\section{信息率失真函数 \texorpdfstring{\(R(D)\)}{R(D)}}
\warmth{0}\quad\whisper{这一节把互信息变成压缩极限：如果允许一点失真，最少需要保留多少信息？}

\begin{strand}
有失真信源编码的核心问题是：在允许平均失真不超过 \(D\) 的条件下，寻找最优的
试验信道 \(\answerin[2.5cm]{p(\hat x|x)}\)，使得互信息
\(\info(X;\hat X)\) 最小。这个最小值就是 \(\answerin[1.5cm]{R(D)}\)。
\end{strand}

\block{从无失真到允许失真}
\trigger{当题目不再要求 \(\hat X=X\)，而是给出一个“平均误差不超过 \(D\)”时，就进入率失真框架。}

在数字通信系统中，信源编码器与译码器的整体作用可以抽象为
\[
X \xrightarrow{\quad\text{编码/译码}\quad} \hat X.
\]
无失真编码要求 \(\hat X=X\)；有失真编码只要求二者之间的
\(\answerin[2cm]{平均失真}\) 不超过某个允许值 \(D\)。

\begin{definition}[有失真编码的目标]
设信源 \(X\) 的字母表为 \(\Xcal=\{x_1,x_2,\cdots,x_n\}\)，重构符号字母表为
\(\hat{\Xcal}=\{\hat x_1,\hat x_2,\cdots,\hat x_m\}\)。
有失真信源编码的目标是：在
\[
\mathbb{E}\left[d(X,\hat X)\right] \leq \answerin[1.2cm]{D}
\]
的条件下，使编码所需的信息率尽可能 \(\answerin[1.5cm]{小}\)。
\end{definition}

\block{失真测度}
\begin{definition}[失真测度]
失真测度是一个非负函数
\[
d:\Xcal\times\hat{\Xcal}\to \mathbb{R}^{+},
\]
其中 \(d(x_i,\hat x_j)\) 表示把 \(x_i\) 重构为 \(\hat x_j\) 所产生的失真程度。
通常要求
\[
d(x_i,\hat x_j)\geq 0,\qquad d(x_i,x_i)=\answerin[1cm]{0}.
\]
\end{definition}

\begin{example}[汉明失真]
对于离散信源，汉明失真只关心“对”或“错”：
\[
d(x_i,\hat x_j)=
\begin{cases}
\answerin[1cm]{0}, & i=j,\\
\answerin[1cm]{1}, & i\neq j.
\end{cases}
\]
它的含义是：译码正确时失真为零，任何错误都计为同样大小的失真。
\end{example}

\begin{example}[平方误差失真]
对于连续信源，常用的失真测度是
\[
d(x,\hat x)=\answerin[2cm]{(x-\hat x)^2}.
\]
若 \(X,\hat X\) 为向量，则可推广为欧氏距离平方
\(\|\mathbf{x}-\hat{\mathbf{x}}\|^2\)。
\end{example}

\begin{yourturn}
汉明失真和平方误差失真的侧重点有何不同？
\[
\text{汉明失真：}\quad d\text{ 只与 }\answerin[2cm]{\text{是否相同}}\text{ 有关。}
\]
\[
\text{平方误差失真：}\quad d\text{ 还与 }\answerin[2cm]{\text{误差大小}}\text{ 有关。}
\]
\workspace[2]
\end{yourturn}


\block{平均失真与允许失真度}
设信源分布为 \(\pmf(x_i)\)，试验信道为 \(\pmf(\hat x_j|x_i)\)，则联合分布为
\[
\pmf(x_i,\hat x_j)=\pmf(x_i)\pmf(\hat x_j|x_i).
\]

\begin{definition}[平均失真]
平均失真定义为
\[
\bar d
=
\sum_{i=1}^{n}\sum_{j=1}^{m}
\pmf(x_i,\hat x_j)d(x_i,\hat x_j)
=
\mathbb{E}\left[d(X,\hat X)\right].
\]
\end{definition}

\begin{definition}[允许失真度]
允许失真度 \(D\) 是平均失真的上界。一个试验信道可行，当且仅当
\[
\bar d
=
\sum_{i=1}^{n}\sum_{j=1}^{m}
\pmf(x_i)\pmf(\hat x_j|x_i)d(x_i,\hat x_j)
\leq \answerin[1cm]{D}.
\]
\end{definition}

\begin{yourturn}
当 \(D\) 变小时，对重构质量的要求变 \(\answerin[1.5cm]{高}\)，
因此所需信息率通常会变 \(\answerin[1.5cm]{大}\)；
当 \(D\) 变大时，信息率可以变 \(\answerin[1.5cm]{小}\)。
\workspace[2]
\end{yourturn}

\block{试验信道}
\begin{definition}[试验信道]
在率失真理论中，信源编码与译码的整体过程抽象为条件概率分布
\(\pmf(\hat x_j|x_i)\)，称为 \textbf{试验信道}。
它不是一个物理信道，而是描述编码/译码造成的随机映射。
\end{definition}

所有满足平均失真约束的试验信道构成集合
\[
\mathcal{P}_D
=
\left\{
\pmf(\hat x|x):
\mathbb{E}\left[d(X,\hat X)\right]\leq D
\right\}.
\]
也就是说，\(\pmf(\hat x|x)\in \mathcal{P}_D\) 当且仅当
\[
\sum_{i=1}^{n}\sum_{j=1}^{m}
\pmf(x_i)\pmf(\hat x_j|x_i)d(x_i,\hat x_j)
\leq D.
\]

\block{互信息作为“必须保留的信息量”}
\begin{proposition}[互信息的概率形式]
原信源 \(X\) 与重构信源 \(\hat X\) 之间的互信息为
\[
\info(X;\hat X)
=
\sum_{i=1}^{n}\sum_{j=1}^{m}
\pmf(x_i)\pmf(\hat x_j|x_i)
\log
\frac{\pmf(\hat x_j|x_i)}{\pmf(\hat x_j)},
\]
其中
\[
\pmf(\hat x_j)=\sum_{i=1}^{n}\pmf(x_i)\pmf(\hat x_j|x_i).
\]
\end{proposition}

互信息衡量：重构信源 \(\hat X\) 中保留了多少关于原信源 \(X\) 的信息。
在有失真压缩中，\(\info(X;\hat X)\) 越小，意味着需要保留的信息越少，
相应编码所需的信息率越 \(\answerin[1cm]{低}\)。


\block{信息率失真函数的定义}
\begin{definition}[信息率失真函数]
信息率失真函数定义为
\[
R(D)
=
\min_{\pmf(\hat x|x)\in \mathcal{P}_D}
\info(X;\hat X).
\]
展开写即
\[
R(D)
=
\min_{\pmf(\hat x|x)}
\info(X;\hat X)
\]
满足约束
\[
\sum_{i=1}^{n}\sum_{j=1}^{m}
\pmf(x_i)\pmf(\hat x_j|x_i)d(x_i,\hat x_j)
\leq D,
\]
\[
\sum_{j=1}^{m}\pmf(\hat x_j|x_i)=1,
\quad i=1,2,\cdots,n,
\]
\[
\pmf(\hat x_j|x_i)\geq 0,
\quad i=1,2,\cdots,n,\ j=1,2,\cdots,m.
\]
\end{definition}

\begin{strand}
\(R(D)\) 的含义是：\textbf{在平均失真不超过 \(D\) 的条件下，编码所需的最小信息率}。
它是一个 \(\answerin[3cm]{\text{带约束的优化问题}}\)，优化变量是试验信道
\(\pmf(\hat x|x)\)，目标函数是互信息。
\end{strand}

\block{率失真函数的基本性质}

\begin{proposition}[非负性]
由于互信息非负，\(\info(X;\hat X)\geq 0\)，所以
\[
R(D)\geq \answerin[1cm]{0}.
\]
\end{proposition}

\begin{proposition}[单调非增性]
若 \(D_1<D_2\)，则
\[
R(D_1)\ \answerin[1.5cm]{\geq}\ R(D_2).
\]
直观含义：允许的失真越大，可选的试验信道越多，所需的最小信息率不会增加。
\end{proposition}

\begin{proposition}[凸性]
率失真函数 \(R(D)\) 是关于 \(D\) 的凸函数，即对 \(0\leq \lambda\leq 1\)，有
\[
R(\lambda D_1+(1-\lambda)D_2)
\leq
\lambda R(D_1)+(1-\lambda)R(D_2).
\]
\end{proposition}

\begin{proof}
设两个最优试验信道分别达到 \((D_1,R(D_1))\) 和 \((D_2,R(D_2))\)。
引入辅助变量 \(\theta\in\{1,2\}\)，以概率 \(\lambda\) 取 \(\theta=1\)、
以概率 \(1-\lambda\) 取 \(\theta=2\)。令 \(\theta\) 与 \(X\) 独立，
并在 \(\theta=k\) 时使用第 \(k\) 个试验信道 \(p_k(\hat x|x)\)。

这样得到的新条件分布为
\[
p(\hat x|x)=\lambda p_1(\hat x|x)+(1-\lambda)p_2(\hat x|x),
\]
其平均失真为
\[
\lambda D_1+(1-\lambda)D_2.
\]

下面估计其互信息。由数据处理不等式，
\[
\info(X;\hat X)\leq \info(X;\hat X,\theta).
\]
而
\[
\info(X;\hat X,\theta)=\info(X;\theta)+\info(X;\hat X|\theta).
\]
因为 \(\theta\) 与 \(X\) 独立，\(\info(X;\theta)=0\)，所以
\[
\info(X;\hat X|\theta)
=\lambda\info(X;\hat X_1)+(1-\lambda)\info(X;\hat X_2)
=\lambda R(D_1)+(1-\lambda)R(D_2).
\]
因此
\[
R(\lambda D_1+(1-\lambda)D_2)
\leq
\lambda R(D_1)+(1-\lambda)R(D_2).
\]
\end{proof}

\begin{yourturn}
凸性的关键动作是：用 \(\answerin[2cm]{\text{时间共享}}\) 把两个编码方案组合起来。
新方案的平均失真是两个方案对应量的 \(\answerin[2cm]{\text{加权平均}}\)，
而新方案的互信息 \(\info(X;\hat X)\) 是加权平均的 \(\answerin[2cm]{\text{上界}}\)。
\workspace[2]
\end{yourturn}

\begin{proposition}[连续性]
在通常条件下，\(R(D)\) 是关于 \(D\) 的连续函数。
\end{proposition}


\block{两个特殊点}

\begin{proposition}[零失真点]
当 \(D=0\) 时，有失真编码退化为无失真编码。

若信源为离散信源，且零失真等价于完全正确恢复，则
\[
R(0)=\info(X;X)=\answerin[2cm]{H(X)}.
\]

若信源为连续信源，在平方误差等常见失真度量下，随着 \(D\to 0\)，
通常有
\[
R(D)\to \answerin[2cm]{\infty}.
\]
\end{proposition}

\begin{proposition}[最大允许失真]
存在一个阈值 \(D_{\max}\)，使得当 \(D\geq D_{\max}\) 时，
可以不传输任何关于 \(X\) 的信息，而在接收端固定输出某个重构符号 \(\hat x\)。
此时 \(X\) 与 \(\hat X\) 独立，\(\info(X;\hat X)=\answerin[1cm]{0}\)，故
\[
R(D)=0,\qquad D\geq D_{\max}.
\]
其中
\[
D_{\max}
=
\min_{\hat x\in \hat{\Xcal}}
\sum_{i=1}^{n}\pmf(x_i)d(x_i,\hat x).
\]
\end{proposition}

\begin{yourturn}
补全 \(D_{\max}\) 的直观含义：在 \(\answerin[2cm]{\text{不发送信息}}\)
的情况下，接收端固定选择一个最优重构符号，
所能达到的最小平均失真为 \(\answerin[2.5cm]{D_{\max}}\)。
\workspace[2]
\end{yourturn}

\block{率失真曲线的整体图像}
综合以上性质，\(R(D)\) 的典型形状满足：
\[
R(0)=H(X),\qquad R(D)\text{ 单调非增且凸},\qquad R(D)=0\ (D\geq D_{\max}).
\]
它刻画了 \(\answerin[3cm]{\text{压缩率}}\) 与
\(\answerin[3cm]{\text{重构质量}}\) 之间的最优折中。

\begin{yourturn}
用一句话概括 \(R(D)\) 回答的问题：
\[
\text{若允许平均失真不超过 }D\text{，最少需要保留多少}
\answerin[2cm]{\text{比特}}\text{ 来描述信源？}
\]
\workspace[2]
\end{yourturn}

\block{本节小结}
\noindent\begin{tabularx}{\linewidth}{@{}l L@{}}
\toprule
{\color{indigo}概念} & \multicolumn{1}{l}{\color{indigo}作用}\\
\midrule
失真测度 \(d(x,\hat x)\) & 刻画原符号与重构符号之间的差异。\\
平均失真 \(\bar d=\mathbb{E}[d(X,\hat X)]\) & 描述整体重构误差。\\
允许失真度 \(D\) & 平均失真的上界。\\
试验信道 \(\pmf(\hat x|x)\) & 描述编码/译码整体造成的随机映射。\\
率失真函数 \(R(D)\) & 在失真约束下，编码所需的最小信息率。\\
\bottomrule
\end{tabularx}

\recall{为什么 \(R(D)\) 关于 \(D\) 是单调非增的？从可行集合
\(\mathcal{P}_D\) 的变化来解释。}

\loose{下一章可以接一个具体信源的 \(R(D)\) 计算，例如二元对称信源在汉明失真下的率失真函数。}
